\documentclass[12pt,a4paper]{article}
\usepackage[margin=25mm, headheight=15pt, headsep=10pt]{geometry}
\usepackage{fontspec}
\usepackage[table]{xcolor} % Note: table option is required for rowcolors
\usepackage{titlesec}
\usepackage{tocloft}
\usepackage{fancyhdr}
\usepackage{booktabs}
\usepackage{tcolorbox}
\tcbuselibrary{skins, breakable}
\usepackage[colorlinks=true, linkcolor=emerald, urlcolor=emerald, bookmarks=true]{hyperref}

\definecolor{emerald}{RGB}{0,102,51}
\definecolor{darkred}{RGB}{139,0,0}
\definecolor{lightgray}{RGB}{245,245,245}

\usepackage[nil]{babel}
\babelprovide[import, main]{bengali}
\babelprovide[import]{arabic}
\babelfont{rm}[Renderer=HarfBuzz,Scale=1.0]{Noto Serif Bengali}
\babelfont{sf}[Renderer=HarfBuzz]{Noto Sans Bengali}
\babelfont{tt}[Renderer=HarfBuzz]{Noto Sans Bengali}
\babelfont[arabic]{rm}[Renderer=HarfBuzz,Scale=1.1]{Noto Naskh Arabic}

% Section styling
\titleformat{\section}{\Large\bfseries\color{emerald}}{\thesection.}{0.5em}{}
\titleformat{\subsection}{\large\bfseries\color{emerald!85!black}}{\thesubsection.}{0.5em}{}
\titleformat{\subsubsection}{\normalsize\bfseries\color{emerald!70!black}}{\thesubsubsection.}{0.5em}{}
\titlespacing*{\section}{0pt}{18pt}{8pt}

% Fancy Header/Footer setup
\pagestyle{fancy}
\fancyhf{} % clear all header and footer fields
\fancyhead[L]{\color{emerald}\small\textbf{\nouppercase{\leftmark}}}
\fancyfoot[L]{\scriptsize\color{black!50}{\lat{AI}}-সংকলিত, আলেম-যাচাইকৃত নয়}
\fancyfoot[C]{\color{emerald!70!black}\thepage}
\renewcommand{\headrulewidth}{0.4pt}
\renewcommand{\headrule}{\hbox to\headwidth{\color{emerald}\leaders\hrule height \headrulewidth\hfill}}
\renewcommand{\footrulewidth}{0pt}

% Custom boxes
% 1. Ayat/Hadith Box
\newtcolorbox{ayatbox}[1][]{
    enhanced, breakable, 
    colback=emerald!5, 
    colframe=emerald, 
    boxrule=0.5pt, 
    arc=3pt, 
    left=10pt, right=10pt, top=10pt, bottom=10pt,
    #1
}

% 2. Quote/Translation Box
\newtcolorbox{quotebox}[1][]{
    enhanced, breakable,
    frame hidden,
    colback=lightgray,
    borderline west={3pt}{0pt}{emerald},
    left=15pt, right=10pt, top=10pt, bottom=10pt,
    sharp corners,
    #1
}

% 3. AI-generated notice box
\newtcolorbox{warnbox}[1][]{
    enhanced, breakable,
    colback=red!4,
    colframe=darkred,
    boxrule=0.8pt,
    arc=3pt,
    left=10pt, right=10pt, top=10pt, bottom=10pt,
    #1
}

% Commands
\newcommand{\arab}[1]{\foreignlanguage{arabic}{#1}}
\newcommand{\kib}[1]{\textcolor{emerald}{\textbf{#1}}}

% Latin (English) fallback: Noto Serif Bengali has NO Latin glyphs
\newfontfamily\latfont[Renderer=HarfBuzz]{Noto Serif}
\newcommand{\lat}[1]{{\latfont #1}}

\setlength{\parskip}{5pt}
\setlength{\parindent}{0pt}

% Fix for missing bullet in Noto Serif Bengali
\renewcommand{\labelitemi}{$\bullet$}



\begin{document}
\begin{center}{\Large\bfseries\color{emerald} \lat{00}-পাঠ-তালিকা}\end{center}
\vspace{1em}
\section*{নামাজের তিলাওয়াতের অর্থ শেখা — পাঠ-তালিকা}

\begin{warnbox}
 \textbf{গুরুত্বপূর্ণ নোটিশ (\lat{Important} \lat{Notice}):} এই কন্টেন্টটি \lat{AI} (কৃত্রিম বুদ্ধিমত্তা) দিয়ে প্রস্তুত ও সংকলিত। \lat{AI} কঠোর সূত্র-মান নীতি মেনে শুধু নির্ভরযোগ্য উৎস থেকে তথ্য সংগ্রহ করেছে — কেবল সহীহ/হাসান হাদিস ও সহীহ সনদের আসার রাখা হয়েছে, দুর্বল সনদ বাদ দেওয়া হয়েছে। তবে এটি কোনো আলেম বা মুহাদ্দিস কর্তৃক যাচাইকৃত নয়।
\end{warnbox}

\begin{quote}
\textbf{উদ্দেশ্য:} নামাজের প্রতিটি অংশে (তাকবীর থেকে সালাম পর্যন্ত) যা কিছু পড়া/বলা হয়, তার প্রতিটি আরবি শব্দের অর্থ জেনে, বুঝে এবং মনে রেখে নামাজ পড়া — যাতে \textbf{\lat{khushu} (খুশু)} অর্জিত হয়।
\end{quote}

\begin{quote}
\textbf{পড়ার ধারা:} প্রতিটি ফাইলে আছে — \textbf{সূত্র কার্ড} $\rightarrow$ \textbf{মূল আরবি} $\rightarrow$ \textbf{বাংলা উচ্চারণ} $\rightarrow$ \textbf{শব্দ-শব্দ অর্থ টেবিল} $\rightarrow$ \textbf{পূর্ণ বাংলা অনুবাদ} $\rightarrow$ \textbf{সংক্ষিপ্ত ব্যাখ্যা (খুশু-কেন্দ্রিক)} $\rightarrow$ \textbf{সংশ্লিষ্ট হাদিস}।
\end{quote}

\medskip\hrule\medskip

\subsection*{কুরআনিক তিলাওয়াত}

\begin{table}[h]\centering\small
\begin{tabular}{llll}
\# & বিষয় & ফাইল & \lat{PDF} \\
\hline
১ & সূরা ফাতিহা — শব্দ-শব্দ অর্থ & `কুরআন/\lat{01}-সূরা-ফাতিহা-শব্দ-অর্থ.\lat{md}` & `কুরআন/\lat{01}-সূরা-ফাতিহা-শব্দ-অর্থ.\lat{pdf}` \\
২ & ছোট সূরা সমূহ (ইখলাস, ফালাক, নাস, কাওসার, আসর, মাউন, কুরাইশ) & `কুরআন/\lat{02}-ছোট-সূরা-সমূহ.\lat{md}` & `কুরআন/\lat{02}-ছোট-সূরা-সমূহ.\lat{pdf}` \\
৩ & আয়াতুল কুরসি ও অন্যান্য & `কুরআন/\lat{03}-আয়াতুল-কুরসি-এবং-অন্যান্য.\lat{md}` & `কুরআন/\lat{03}-আয়াতুল-কুরসি-এবং-অন্যান্য.\lat{pdf}` \\
\hline
\end{tabular}
\end{table}

\subsection*{হাদিস — খুশু ও তিলাওয়াতের ফজিলত}

\begin{table}[h]\centering\small
\begin{tabular}{llll}
\# & বিষয় & ফাইল & \lat{PDF} \\
\hline
৪ & খুশুর গুরুত্ব — সহীহ হাদিস & `হাদিস/\lat{01}-খুশুর-গুরুত্ব-হাদিস.\lat{md}` & `হাদিস/\lat{01}-খুশুর-গুরুত্ব-হাদিস.\lat{pdf}` \\
৫ & তিলাওয়াতের ফজিলত — সহীহ হাদিস & `হাদিস/\lat{02}-তিলাওয়াতের-ফজিলত-হাদিস.\lat{md}` & `হাদিস/\lat{02}-তিলাওয়াতের-ফজিলত-হাদিস.\lat{pdf}` \\
\hline
\end{tabular}
\end{table}

\subsection*{নামাজের প্রতিটি অংশের তিলাওয়াত ও অর্থ (ধাপে ধাপে)}

\begin{table}[h]\centering\small
\begin{tabular}{llll}
\# & অংশ & ফাইল & \lat{PDF} \\
\hline
৬ & তাকবীরে তাহরিমা — \textit{\lat{Allahu} \lat{Akbar}} & `\lat{01}-তাকবীরে-তাহরিমা.\lat{md}` & `\lat{01}-তাকবীরে-তাহরিমা.\lat{pdf}` \\
৭ & সানা — \textit{\lat{Subhanaka} \lat{Allahumma}...} & `\lat{02}-সানা-ইস্তিফতাহ.\lat{md}` & `\lat{02}-সানা-ইস্তিফতাহ.\lat{pdf}` \\
৮ & তাউয়্যুয় ও বিসমিল্লাহ — \textit{\lat{A'udhu} + \lat{Bismillah}} & `\lat{03}-তাউয়্যুয-বিসমিল্লাহ.\lat{md}` & `\lat{03}-তাউয়্যুয-বিসমিল্লাহ.\lat{pdf}` \\
৯ & সূরা ফাতিহা — অর্থসহ & `\lat{04}-সূরা-ফাতিহা-অর্থ.\lat{md}` & `\lat{04}-সূরা-ফাতিহা-অর্থ.\lat{pdf}` \\
১০ & সূরা মিলানো — কিরাআত & `\lat{05}-সূরা-মিলানো-কিরাআত.\lat{md}` & `\lat{05}-সূরা-মিলানো-কিরাআত.\lat{pdf}` \\
১১ & রুকু — তাসবিহ ও দোয়া & `\lat{06}-রুকু-তাসবিহ-দোয়া.\lat{md}` & `\lat{06}-রুকু-তাসবিহ-দোয়া.\lat{pdf}` \\
১২ & কাউমা — রুকু থেকে দাঁড়ানোর দোয়া & `\lat{07}-কাউমা-দোয়া.\lat{md}` & `\lat{07}-কাউমা-দোয়া.\lat{pdf}` \\
১৩ & সিজদা — তাসবিহ ও দোয়া & `\lat{08}-সিজদা-তাসবিহ-দোয়া.\lat{md}` & `\lat{08}-সিজদা-তাসবিহ-দোয়া.\lat{pdf}` \\
১৪ & জালসা — দুই সিজদার মাঝের বৈঠক & `\lat{09}-জালসা-দোয়া.\lat{md}` & `\lat{09}-জালসা-দোয়া.\lat{pdf}` \\
১৫ & তাশাহহুদ — প্রথম বৈঠক & `\lat{10}-তাশাহহুদ-প্রথম.\lat{md}` & `\lat{10}-তাশাহহুদ-প্রথম.\lat{pdf}` \\
১৬ & তাশাহহুদ ও দরুদ ইবরাহিম — শেষ বৈঠক & `\lat{11}-তাশাহহুদ-শেষ-দরুদ.\lat{md}` & `\lat{11}-তাশাহহুদ-শেষ-দরুদ.\lat{pdf}` \\
১৭ & শেষ দোয়া ও সালাম & `\lat{12}-শেষ-দোয়া-সালাম.\lat{md}` & `\lat{12}-শেষ-দোয়া-সালাম.\lat{pdf}` \\
১৮ & নামাজের পরের যিকর ও দোয়া & `\lat{13}-নামাজের-পর-যিকর-দোয়া.\lat{md}` & `\lat{13}-নামাজের-পর-যিকর-দোয়া.\lat{pdf}` \\
১৯ & কুনূত, নফল ও বিশেষ দোয়া & `\lat{14}-কুনূত-নফল-বিশেষ-দোয়া.\lat{md}` & `\lat{14}-কুনূত-নফল-বিশেষ-দোয়া.\lat{pdf}` \\
২০ & ব্যবহারিক গাইড — খুশু জনন & `\lat{15}-ব্যবহারিক-গাইড-খুশু-জনন.\lat{md}` & `\lat{15}-ব্যবহারিক-গাইড-খুশু-জনন.\lat{pdf}` \\
\hline
\end{tabular}
\end{table}

\subsection*{সাহাবিদের ঘটনা (খুশু ও নামাজ সম্পর্কিত)}

\begin{table}[h]\centering\small
\begin{tabular}{lll}
\# & ঘটনা & ফাইল \\
\hline
২১ & নবীজি \arab{ﷺ} — রাতের নামাজে পা ফুলে যাওয়া, 'কৃতজ্ঞ বান্দা' & `সাহাবিদের-ঘটনা/\lat{01}-নবীজির-খুশু-নামাজ.\lat{md}` \\
২২ & হুযাইফা (রাঃ) — নবীজির পেছনে রাতের নামাজ & `সাহাবিদের-ঘটনা/\lat{02}-হুযাইফার-রাতের-নামাজ.\lat{md}` \\
২৩ & সাহাবীদের নামাজে কান্না — হৃদয়ের কোমলতা & `সাহাবিদের-ঘটনা/\lat{03}-সাহাবীদের-নামাজে-কান্না.\lat{md}` \\
\hline
\end{tabular}
\end{table}

\medskip\hrule\medskip

\subsection*{মূল হাদিস (স্মরণীয় উদ্ধৃতি)}

\begin{quote}
\textbf{\lat{“}নামাজ পড়ো, যেমন আমাকে নামাজ পড়তে দেখেছ।\lat{”}} — সহীহ বুখারি ৬৩১
\end{quote}

\begin{quote}
\textbf{\lat{“}যে ব্যক্তি সূরা ফাতিহা পড়েনি, তার নামাজ নেই।\lat{”}} — সহীহ বুখারি ৭৫৬
\end{quote}

\begin{quote}
\textbf{\lat{“}বান্দা তার রবের সবচেয়ে নিকটবর্তী হয় সিজদার অবস্থায় — তাই সিজদায় বেশি দোয়া করো।\lat{”}} — সহীহ মুসলিম ৪৮২
\end{quote}

\begin{quote}
\textbf{\lat{“}যখন ইমাম 'গাইরিল মাগদুবি আলাইহিম ওয়ালাদ্দাল্লিন' বলেন, তোমরা আমিন বলো; কারণ ফেরেশতারাও আমিন বলে।\lat{”}} — সহীহ বুখারি ৭৮০
\end{quote}

\medskip\hrule\medskip

\subsection*{গঠন পদ্ধতি}

\begin{itemize}
\item \textbf{মার্কডাউন (.\lat{md}):} দ্রুত পড়া ও এডিটের জন্য
\item \textbf{পিডিএফ (.\lat{pdf}):} আরামে পড়ার জন্য, লুয়াটেক + নোটো বাংলা/আরবি ফন্ট দিয়ে তৈরি
\item প্রতিটি সাহাবির ঘটনা আলাদা ফাইল — প্রতিটি ঘটনা স্বতন্ত্র পাঠযোগ্য
\item `\lat{shared}/\lat{preamble}.\lat{tex}` — \lat{PDF}-এর সাধারণ ফন্ট ও স্টাইল; `\lat{scripts}/\lat{md2tex}.\lat{py}` — \lat{md} $\rightarrow$ \lat{tex} রূপান্তর; `\lat{Makefile}` — `\lat{make} \lat{build}\_\lat{pdfs}` দিয়ে সব \lat{PDF}
\end{itemize}
\medskip\hrule\medskip

\textit{শেষ আপডেট: আগস্ট ২০২৬}

\end{document}
